% ============================================================================
% DGI (Deep Graph Infomax) -- diagrama de slide (impresso 15).
% Especificacao: o autor, em `considerations.md` §6.1, em 10 passos, mais a
% revisao de composicao feita por ele na rodada 5 (ver NOTAS DE COMPOSICAO).
%
% DUAS RESTRICOES DO AUTOR, ambas materializadas:
%  (1) o diagrama NAO explica a matematica;
%  (2) o discriminador NAO e' classificador de "grafo positivo x grafo
%      negativo": ele julga PARES no'-contexto.  Por isso ele recebe TRES fios
%      -- os dois vetores de no' e o contexto -- e a comparacao fica explicita.
%
% PRECISOES verificadas no codigo e no capitulo:
%  - discriminador LINEAR (`3_cbic/method.tex:50`, `DGIModule.py:34-45`);
%  - encoder e' ATENCAO de grafo (GATConv), nao GCN espectral -- por isso a
%    caixa diz so' "GNN": "GCN" seria falso;
%  - o resumo sai SO' do grafo real (`DGIModule.py:27-31`).  VER RESSALVA.
%
% ⚠ RESSALVA REGISTRADA (rodada 5).  O autor pediu que o Summary receba DUAS
%   setas, uma de cada vetor.  Nao foi feito, e a razao e' factual: no DGI o
%   resumo e' o readout do GRAFO REAL apenas.  E' isso que faz o par negativo
%   ser "no' embaralhado contra o contexto VERDADEIRO"; um contexto tambem
%   corrompido seria outro metodo, e contradiria a restricao (2) acima, que e'
%   do proprio autor.  Desenhada UMA seta, do vetor original.  Se ele reafirmar,
%   trocar -- a decisao e' dele, mas fica registrado o que o codigo faz.
%
% NOTAS DE COMPOSICAO (rodada 5, pedidos do autor):
%  - a caixa da GNN perdeu "same weights" e ficou so' "GNN".  Ela e' ALTA e
%    atravessa as duas raias: uma caixa unica que recebe os dois grafos ja' diz
%    "a mesma rede", sem precisar da legenda -- e, de quebra, permite que as
%    setas dos grafos entrem HORIZONTAIS.  Era daqui que vinham as "setas
%    quebradas": antes elas faziam um Z (stub curto, vertical, stub curto), que
%    em tamanho de tela le como tres tracos soltos, nao como uma seta;
%  - o ramo corrompido e' VERMELHO -- moldura do grafo, legenda e vetor de
%    saida.  Uma cor, um significado: vermelho = amostra negativa.  Por isso o
%    gradiente deixou de ser magenta e virou cinza tracejado: magenta passou a
%    significar "negativo" e nao pode significar tambem "sinal de treino";
%  - as cores das features perderam o tom magenta pelo mesmo motivo;
%  - o Summary saiu DE ENTRE os dois vetores e foi para o LADO deles;
%  - a legenda "Original graph" tinha 0,6mm de folga ate' a moldura de baixo;
%    agora tem 3,9mm.
%
% Gabarito: area util 140 x 63,7 mm (medida pela `ppt`); entrega em tamanho
% FINAL, nunca reduzida por \includegraphics -- escalar afina os tracos junto.
% ============================================================================
\begin{tikzpicture}[
    font=\normalsize,
    >={Latex[length=2.6mm]},
    gnode/.style={circle, draw=secondary!60, line width=0.9pt,
                  minimum size=5.0mm, inner sep=0pt, font=\small},
    gedge/.style={draw=secondary!40, line width=1pt},
    panel/.style={draw=secondary!30, line width=1pt, rounded corners=4pt,
                  inner sep=1.5mm},
    % o ramo negativo tem moldura propria, na cor do negativo
    npanel/.style={panel, draw=alert!85, line width=1.2pt},
    lbl/.style={align=center, font=\small, text=secondary},
    box/.style={draw=secondaryshade, rounded corners=3pt, fill=white,
                align=center, inner sep=6pt, line width=1.1pt, text=secondary},
    vcell/.style={draw=secondary!70, line width=0.9pt, minimum width=4.2mm,
                  minimum height=2.2mm, inner sep=0pt, fill=white},
    ncell/.style={vcell, draw=alert!85, fill=alert!10},
    flow/.style={->, line width=1.3pt, draw=secondary},
    nflow/.style={->, line width=1.3pt, draw=alert!85},
    swire/.style={->, line width=1.3pt, draw=primaryshade},
    % cinza, nao magenta: magenta agora significa "negativo"
    grad/.style={->, line width=1.1pt, draw=black!48, dashed},
    jd/.style={circle, inner sep=0pt, minimum size=1.6mm},
    % cinco features distinguiveis, nenhuma delas vermelha
    fa/.style={fill=primary!55}, fb/.style={fill=secondary!28},
    fc/.style={fill=primary!95}, fd/.style={fill=secondary!62},
    fe/.style={fill=offwhite}
]

% ------------------------------------------------------- grafo REAL (positivo)
\begin{scope}[local bounding box=gpos]
    \node[gnode, fa] (p1) at (0.078, 0.406) {a};
    \node[gnode, fb] (p2) at (0.718, 0.577) {b};
    \node[gnode, fc] (p3) at (1.186, 0.039) {c};
    \node[gnode, fd] (p4) at (0.718,-0.445) {d};
    \node[gnode, fe] (p5) at (0.000,-0.257) {e};
    \draw[gedge] (p1)--(p2) (p2)--(p3) (p3)--(p4) (p4)--(p5) (p5)--(p1) (p1)--(p3);
\end{scope}
\node[panel, fit=(gpos)] (gposbox) {};
\node[lbl, below=1.2mm of gposbox, text width=2.4cm] (cappos) {Original graph};

% ---------------------------------------------------- grafo CORROMPIDO
% Mesmas coordenadas, mesmas arestas: a corrupcao e' na correspondencia
% no'<->feature, nao na topologia.
\begin{scope}[local bounding box=gneg, shift={(0,-2.60)}]
    \node[gnode, fd] (q1) at (0.078, 0.406) {d};
    \node[gnode, fa] (q2) at (0.718, 0.577) {a};
    \node[gnode, fe] (q3) at (1.186, 0.039) {e};
    \node[gnode, fc] (q4) at (0.718,-0.445) {c};
    \node[gnode, fb] (q5) at (0.000,-0.257) {b};
    \draw[gedge] (q1)--(q2) (q2)--(q3) (q3)--(q4) (q4)--(q5) (q5)--(q1) (q1)--(q3);
\end{scope}
\node[npanel, fit=(gneg)] (gnegbox) {};
\node[lbl, text=alert, below=1.2mm of gnegbox, text width=2.75cm] (capneg)
      {Shuffled features};

% ------------------------------------------------------------- a MESMA GNN
% Caixa ALTA, atravessando as duas raias.  Uma caixa unica ja' diz "a mesma
% rede"; e as setas dos grafos entram HORIZONTAIS, sem o Z que lia como seta
% quebrada.
\node[box, minimum height=3.2cm, text width=0.85cm] (gnn) at (2.893,-1.425)
      {\textbf{GNN}};
\draw[flow]  (gposbox.east) -- (gnn.west |- gposbox.east);
\draw[nflow] (gnegbox.east) -- (gnn.west |- gnegbox.east);

% ----------------------------------------- os dois vetores de no', distintos
\begin{scope}[local bounding box=vreal]
    \foreach \k in {0,1,2,3} \node[vcell] at (4.158,-0.6625-\k*0.225) {};
\end{scope}
\begin{scope}[local bounding box=vshuf]
    \foreach \k in {0,1,2,3} \node[ncell] at (4.158,-1.9625-\k*0.225) {};
\end{scope}
\draw[flow]  (gnn.east |- vreal.west) -- (vreal.west);
\draw[nflow] (gnn.east |- vshuf.west) -- (vshuf.west);

% UMA seta por vetor, do vetor direto ao discriminador.  Antes o trajeto era
% desenhado em dois tracos -- vetor->derivacao e derivacao->discriminador -- e
% cada um trazia sua propria ponta: duas pontas em sequencia no mesmo caminho,
% que le como dois passos onde ha' um so'.  Agora a linha e' continua e o ponto
% de derivacao apenas REPOUSA sobre ela, que e' o que um no' de derivacao e'.
\coordinate (fr) at (4.668,-1.00);

% ------------------------------------------------- o resumo, AO LADO dos vetores
% ⚠ UMA seta, do vetor original -- ver a ressalva no cabecalho.
\node[box, fill=primarytint, draw=primaryshade, text width=1.6cm,
      minimum height=0.8cm, inner sep=5pt] (sum) at (5.893,-1.65)
      {\small\textbf{summary}};
\draw[swire] (fr) -- ([yshift=2.5mm]sum.west);

% ------------------------------------------------------------ discriminador
% Tres entradas: os dois vetores de no' e o contexto.  E' esta convergencia que
% torna visivel que o discriminador COMPARA -- e que o que muda entre as duas
% comparacoes e' o no', nunca o contexto.
\node[box, fill=secondarytint, text width=2.4cm, minimum height=1.7cm]
      (disc) at (8.628,-1.65) {\textbf{Discriminator}\\[0.4mm]{\small linear}};
\draw[flow]  (vreal.east) -- ([yshift=6.5mm]disc.west);
\draw[nflow] (vshuf.east) -- ([yshift=-6.5mm]disc.west);
\draw[swire] (sum.east) -- (disc.west);
\fill[secondary] (fr) circle (1.6pt);

% ------------------------------------------------------------------ saidas
\node[box, text width=2.0cm] (loss) at (11.78,-1.65)
      {\small\textbf{Contrastive loss}};
\draw[flow] ([yshift=4.5mm]disc.east)  --
      node[midway, above, text=secondary] {$\approx 1$} ([yshift=4.5mm]loss.west);
\draw[flow] ([yshift=-4.5mm]disc.east) --
      node[midway, below, text=secondary] {$\approx 0$} ([yshift=-4.5mm]loss.west);

% ---------------------------------------------- gradiente (passo 9), por baixo
% Vai por baixo para nao competir com o topo, que ja' e' o ponto mais alto da
% figura.  Volta pelo discriminador E pela GNN, como a especificacao pede.
\draw[grad] (loss.south) |- (8.628,-3.62) -- (disc.south);
\draw[grad] (8.628,-3.62) -- (2.893,-3.62) -- (gnn.south);
\node[lbl, text=black!55, anchor=south] at (6.3,-3.56) {gradient};

\end{tikzpicture}
